\section*{Capítulo \ref{ch:b1:ocean-circulation} --- A circulação do oceano}
% values: calc/ch29.py

\begin{solution}{exo:b1:ocean-circulation:1}
(a)~A \qty{4}{m/s} (primeira fórmula): $1000\,C_D = 0.29 + 0.775 + 0.481 = 1.546$,
$T = 1.3 \times \num{1.546e-3} \times 16 = \qty{0.0322}{Pa}$. A \qty{7}{m/s}
(segunda fórmula): $1000\,C_D = 0.60 + 0.497 = 1.097$, $T = \qty{0.0699}{Pa}$. A
\qty{15}{m/s}: $1000\,C_D = 1.665$, $T = \qty{0.487}{Pa}$. (b)~$0.487/0.0699 =
6.97$, contra $(15/7)^2 = 4.59$: o próprio coeficiente de arrasto cresce com o vento,
porque um vento mais forte levanta ondas mais rugosas que agarram melhor o ar.
\end{solution}

\begin{solution}{exo:b1:ocean-circulation:2}
(a)~Hemisfério Norte, \ang{90} à direita de um vento soprando para leste: para o
sul. (b)~O vento sopra para sudoeste; à sua direita está o noroeste.
(c)~Hemisfério Sul, \ang{90} à esquerda de um vento soprando para leste: para o
norte. (d)~O vento sopra para o norte; à sua esquerda está o oeste, que é para o mar
em uma costa oeste. No caso (d) o transporte aponta para o mar: a água da superfície é
reposta por baixo, e água fria ressurge ao longo da costa.
\end{solution}

\begin{solution}{exo:b1:ocean-circulation:3}
(a)~$31.7 + 3.5 - 17.9 = \qty{17.3}{\sverdrup}$, contra o publicado
\qty{17.1}{\sverdrup}: as médias das três parcelas somam o revolvimento médio dentro
do arredondamento e das incertezas ($\pm\qty{3.2}{\sverdrup}$).
(b)~Acima de \qty{1100}{m}: $16.6 + 0.4 = \qty{17.0}{\sverdrup}$ para o norte; abaixo:
$12.0 + 5.9 = \qty{17.9}{\sverdrup}$ para o sul. Os dois quase se equilibram, como devem
fazer quando nenhuma água se acumula ao norte da seção; a pequena diferença está dentro das
incertezas. (c)~$\num{17.1e6} \times \num{3.156e7} = \qty{5.40e14}{m^3}$, ou seja
\qty{5.40e5}{km^3} por ano; $\num{1.3324e9}/(\num{5.40e5}) \approx 2470$ anos.
\end{solution}

\begin{solution}{exo:b1:ocean-circulation:4}
$D_E = 7.6 \times 10/\sqrt{\sin\varphi}$ e $V_0 = 0.0127 \times
10/\sqrt{\sin\varphi}$. A \ang{15}: \qty{149}{m} e \qty{0.250}{m/s}
(\qty{2.5}{\percent} do vento); a \ang{45}: \qty{90}{m} e \qty{0.151}{m/s}
(\qty{1.5}{\percent}); a \ang{60}: \qty{82}{m} e \qty{0.136}{m/s}
(\qty{1.4}{\percent}). Perto do equador $\sin\varphi \to 0$ e as duas fórmulas
divergem; mais fundamentalmente $f \to 0$, de modo que a força de Coriolis necessária para equilibrar
o atrito se anula, e o tempo $1/f$ que o equilíbrio leva para se estabelecer torna-se maior
que o tempo em que os ventos permanecem constantes. Ekman restringiu sua fórmula a $|\varphi| \geq
\ang{10}$.
\end{solution}

\begin{solution}{exo:b1:ocean-circulation:5}
(a)~$f = 2\Omega\sin\ang{11} = \qty{2.783e-5}{s^{-1}}$; $M = 0.07/(\num{2.783e-5}) =
\qty{2515}{kg.m^{-1}.s^{-1}}$, \ang{90} à direita de um vento soprando para oeste:
para o norte. (b)~$Q = (2515/1027) \times \num{5e6} = \qty{1.22e7}{m^3/s}$,
\qty{12.2}{\sverdrup}. (c)~A estimativa fica dentro da incerteza da medida direta
($12.0 \pm 5.5$) e próxima do valor dos ventos de longo prazo
($13.5 \pm 0.3$): o teorema explica o transporte observado.
\end{solution}

\begin{solution}{exo:b1:ocean-circulation:6}
(a)~$f(\ang{2}) = \qty{5.09e-6}{s^{-1}}$; $M = 0.05/(\num{5.09e-6}) =
\qty{9824}{kg.m^{-1}.s^{-1}}$, \qty{9.57}{m^2/s}: a \ang{2} N para o norte
(à direita de um vento de oeste), a \ang{2} S para o sul (à esquerda). (b)~Sobre uma faixa
\qty{1}{m} de comprimento em \omterm{def:b1:spherical-astronomy:geographic}{longitude}, $2 \times 9.57 = \qty{19.1}{m^3/s}$ saem através das
duas \omterm{def:b1:spherical-astronomy:geographic}{latitudes}, a \qty{444.8}{km} uma da outra: $w = 19.1/(\num{4.448e5}) = \qty{4.30e-5}{m/s}$,
cerca de \qty{3.7}{m} por dia. (c)~No equador $f = 0$ e $|\vec M| = T/|f|$ é
infinito: o equilíbrio de Ekman permanente não existe ali, e os números acima, muito
sensíveis à \omterm{def:b1:spherical-astronomy:geographic}{latitude} escolhida, são apenas uma ordem de grandeza. O que sobrevive é
o sinal: de um lado e do outro o transporte aponta para longe do equador, de modo que a
superfície diverge e água fria sobe ao longo dela.
\end{solution}

\begin{solution}{exo:b1:ocean-circulation:7}
(a)~$f(\ang{15}) = \qty{3.775e-5}{s^{-1}}$: $M = 0.06/(\num{3.775e-5}) =
\qty{1590}{kg.m^{-1}.s^{-1}}$ para o norte; $f(\ang{40}) =
\qty{9.375e-5}{s^{-1}}$: $M = \qty{1280}{kg.m^{-1}.s^{-1}}$ para o sul.
(b)~Entre as duas \omterm{def:b1:spherical-astronomy:geographic}{latitudes}, separadas por \qty{2780}{km}, $\partial M_y/\partial y =
-(1590 + 1280)/(\num{2.78e6}) = \qty{-1.03e-3}{kg.m^{-2}.s^{-1}}$: pela
\cref{prop:b1:ocean-circulation:pumping}, $w_E = \qty{-1.005e-6}{m/s}$, uma velocidade
descendente de \qty{31.7}{m} por ano. (c)~$\num{1.005e-6} \times \num{2.78e6} \times
\num{5e6} = \qty{1.40e7}{m^3/s}$, \qty{14.0}{\sverdrup} bombeados para baixo.
\end{solution}

\begin{solution}{exo:b1:ocean-circulation:8}
(a)~$H = 1027 \times 3992 \times \num{17e6} \times 12 = \qty{8.36e14}{W}$,
\qty{0.84}{PW}. (b)~São \qty{70}{\percent} dos \qty{1.2}{PW} medidos: a parte superior real
é mais quente perto da superfície que a média de \qty{15}{\celsius}, e o
\omterm{def:b1:ocean-circulation:gyre}{giro} também transporta calor. (c)~$1027 \times 3992 \times \num{17e6} \times 3 =
\qty{2.1e14}{W}$, \qty{0.21}{PW}, um quarto de (a). Só a água que se torna fria
antes de retornar carrega muito calor: o transporte de calor do Atlântico depende da
formação de água profunda.
\end{solution}

\begin{solution}{exo:b1:ocean-circulation:9}
Sejam $x$, $y$ e $1 - x - y$ as frações de AAIW, NADW e AABW. Balanços de sal e de calor
medidos a partir da AABW:
\[
  -0.52\,x + 0.33\,y = 34.731 - 34.67 = 0.061,\qquad
  2.9\,x + 3.6\,y = 2.78 - 0.4 = 2.38 .
\]
O determinante é $-0.52 \times 3.6 - 0.33 \times 2.9 = -2.829$, e a regra de Cramer
dá $x = (0.061 \times 3.6 - 0.33 \times 2.38)/(-2.829) = 0.20$ e $y = (-0.52
\times 2.38 - 2.9 \times 0.061)/(-2.829) = 0.50$: a amostra é \qty{20}{\percent}
AAIW, \qty{50}{\percent} NADW e \qty{30}{\percent} AABW, um ponto dentro do
triângulo das três águas.
\end{solution}

\begin{solution}{exo:b1:ocean-circulation:10}
(a)~Com $W = u + iv = V_0e^{i\pi/4}e^{(1+i)az}$, cujas partes real e imaginária são
as fórmulas da proposição, $W'' = (1+i)^2a^2W = 2ia^2W$. As duas equações
combinam-se em $A_z(u'' + iv'') = -fv + ifu = if(u + iv)$, ou seja $A_zW'' = ifW$,
o que vale pois $2ia^2A_z = if$ para $a^2 = f/(2A_z)$. (b)~$W'(0) =
(1+i)aV_0e^{i\pi/4} = \sqrt2\,aV_0\,e^{i\pi/2} = i\sqrt2\,aV_0$: a tensão na superfície
$\rho A_zW'(0)$ aponta para o norte, com módulo $\rho A_z\sqrt2\,aV_0 = \rho A_z\sqrt2
\sqrt{f/(2A_z)}\;T/(\rho\sqrt{fA_z}) = T$. (c)~$\int_{-\infty}^0 W\,\mathrm dz =
V_0e^{i\pi/4}/((1+i)a) = V_0/(\sqrt2\,a)$, um número real: o transporte não tem
parte para o norte, e sua parte para leste é $\rho V_0/(\sqrt2\,a) =
T/\sqrt{fA_z}\cdot\sqrt{A_z/f} = T/f$. Este é o teorema: módulo $T/f$, \ang{90} à
direita de um vento soprando para o norte.
\end{solution}

\begin{solution}{exo:b1:ocean-circulation:11}
(a)~Um passo para o norte $\mathrm dy = a_\oplus\,\mathrm d\varphi$ muda $f = 2\Omega
\sin\varphi$ em $2\Omega\cos\varphi\,\mathrm d\varphi$, logo $\beta = 2\Omega
\cos\varphi/a_\oplus = 2 \times \num{7.292e-5} \times \cos\ang{30}/(\num{6.371e6}) =
\qty{1.98e-11}{m^{-1}.s^{-1}}$. (b)~$R/\beta = \num{e-6}/(\num{1.98e-11}) =
\qty{5.04e4}{m}$, cerca de \qty{50}{km}. (c)~Interior: $\num{30e6}/(\num{5e6} \times
1000) = \qty{6.0e-3}{m/s}$, \qty{6}{mm/s}; corrente de contorno: $\num{30e6}/(\num{5.04e4}
\times 1000) = \qty{0.60}{m/s}$, cem vezes mais rápida, a razão das larguras. A
Corrente do Golfo, a \qty{1.8}{m/s}, é ainda mais rápida: a corrente real é mais estreita
do que esta estimativa, seu escoamento se concentra perto da superfície, e sua recirculação soma-se
a seu transporte.
\end{solution}

\begin{solution}{exo:b1:ocean-circulation:12}
(a)~Onde o mar é mais profundo que \qty{1000}{m}, a coluna acima de \qty{1000}{m}
tem no máximo \qty{1}{km} de água, e em outros lugares menos: o volume acima de
\qty{1000}{m} é no máximo $\num{361.84e6} \times 1 = \qty{3.618e8}{km^3}$, de modo que o
volume abaixo é de pelo menos $\num{1.3324e9} - \num{3.618e8} = \qty{9.71e8}{km^3}$.
(b)~$\num{26e6} \times \num{3.156e7} = \qty{8.21e14}{m^3}$ por ano,
\qty{8.21e5}{km^3}; $\num{9.71e8}/(\num{8.21e5}) \approx 1200$ anos, da mesma ordem que
os mil anos da correia transportadora. (c)~$(1/2)^{25/12.3} = 0.244$, cerca de um quarto;
$\num{e-3} \times 25 \times \num{3.156e7} = \qty{7.9e5}{m}$, cerca de \qty{790}{km}.
Após 25 anos o trítio se espalhou não mais do que uma deriva lenta o levaria,
enquanto um quarto dele ainda resta para ser detectado.
\end{solution}

\begin{solution}{pb:b1:ocean-circulation:1}
\begin{enumerate}
  \item Um vento do norte sopra para o sul, ao longo da costa.
  \item $1000\,C_D = 0.60 + 0.071 \times 8 = 1.168$: $C_D = \num{1.168e-3}$.
  \item $T = 1.3 \times \num{1.168e-3} \times 64 = \qty{0.0972}{Pa}$.
  \item $C_D = \num{1.452e-3}$ e $T = \qty{0.272}{Pa}$: a tensão cresce $2.80$ vezes,
        mais que $(12/8)^2 = 2.25$, porque o próprio $C_D$ cresce com o vento.
  \item $\Omega = \qty{7.292e-5}{rad.s^{-1}}$; $f = 2\Omega\sin\ang{38} =
        \qty{8.979e-5}{s^{-1}}$.
  \item $2\pi/f = \qty{7.00e4}{s}$, \qty{19.4}{h}. A força de Coriolis faz girar uma
        corrente por um círculo completo nesse tempo; antes que o vento tenha soprado por
        cerca desse tempo, a água oscila em vez de se estabelecer no
        equilíbrio permanente do teorema.
  \item \ang{90} à direita de um vento soprando para o sul: para oeste, o mar
        aberto.
  \item $M = T/f = 0.0972/(\num{8.979e-5}) = \qty{1082}{kg.m^{-1}.s^{-1}}$.
  \item $1082/1027 = \qty{1.054}{m^2/s}$ por metro de costa.
  \item $D_E = 7.6 \times 8/\sqrt{\sin\ang{38}} = \qty{77.5}{m}$.
  \item $1.054/77.5 = \qty{0.0136}{m/s}$, \qty{1.36}{cm/s}.
  \item $V_0 = 0.0127 \times 8/\sqrt{\sin\ang{38}} = \qty{0.129}{m/s}$, \ang{45} à
        direita do vento: para sudoeste. Os derivadores dão cerca da metade
        do valor de Ekman, uns \qty{0.07}{m/s}.
  \item A água que deixa a costa na \omterm{def:b1:ocean-circulation:ekman}{camada de Ekman} não pode ser reposta através da
        costa, nem por cima: ela é reposta por água que sobe de baixo da camada
        perto da costa, como na seção transversal desenhada em
        \cref{sec:b1:ocean-circulation:upwelling}; as isotermas se curvam para cima em direção à
        costa.
  \item $1.054 \times 1000 \times 86400 = \qty{9.11e7}{m^3}$ por dia por quilômetro de
        costa.
  \item $w = 1.054/(\num{2e4}) = \qty{5.27e-5}{m/s}$, \qty{4.55}{m} por dia.
  \item $77.5/4.55 = 17$ dias.
  \item O ar é resfriado por baixo pela água fria: a camada perto da superfície
        torna-se fria e estável, a umidade pode condensar em neblina e estratos baixos,
        e sem convecção há pouca chuva na costa.
  \item A água ressurgida é rica em nutrientes; eles alimentam o fitoplâncton, que
        alimenta o zooplâncton e os peixes.
  \item Um vento soprando para o norte impele um transporte a \ang{90} à sua direita, para o
        leste: a água da superfície se acumula contra a costa e afunda (\omterm{def:b1:ocean-circulation:upwelling}{afundamento}); a
        \omterm{def:b1:ocean-circulation:upwelling}{ressurgência} cessa e água superficial mais quente retorna à costa.
  \item $f(\ang{20}) = \qty{4.988e-5}{s^{-1}}$: $M = \qty{1948}{kg.m^{-1}.s^{-1}}$,
        $\sin\ang{38}/\sin\ang{20} = 1.80$ vezes mais.
  \item No Hemisfério Sul o transporte está a \ang{90} à esquerda do
        vento. Para empurrá-lo para oeste, para o mar, o vento deve soprar para o norte, a partir do
        sul: nos dois casos o vento sopra ao longo da costa em direção ao equador.
  \item $1000\,C_D = 0.29 + 0.62 + 0.308 = 1.218$; $T = \qty{0.0396}{Pa}$; $M =
        \qty{441}{kg.m^{-1}.s^{-1}}$; a razão é $0.407$ para ambos.
  \item Por metro de costa, $M c_p\,\Delta T = 1082 \times 3992 \times 6 =
        \qty{2.59e7}{W}$, \qty{25.9}{MW}; ao longo de \qty{1000}{km}, \qty{2.59e13}{W},
        \qty{25.9}{TW}.
  \item $\num{31.7e6}/1.054 = \qty{3.01e7}{m}$, cerca de \qty{30000}{km}: \omterm{def:b1:ocean-circulation:upwelling}{ressurgência
        costeira} move muito menos água que uma \omterm{def:b1:ocean-circulation:wbc}{corrente de contorno oeste}.
  \item $Q = 1.054 \times \num{e6} = \qty{1.05e6}{m^3/s}$: o transporte de \omterm{def:b1:ocean-circulation:upwelling}{ressurgência}
        ao longo de \qty{1000}{km} de costa é de \qty{1.05}{\sverdrup}.
\end{enumerate}
\end{solution}
